% Source: ACM_Templates.pdf, PDF pp. 14-16.
% Corrections and validation: docs/validation.md.
\subsection{素数筛与质因数奇偶哈希}
\label{algo:sieve}
\AlgorithmInfo
    {预处理素数、最小质因子、欧拉函数与莫比乌斯函数；可扩展为质因数指数奇偶性的随机指纹。}
    {线性筛时间、空间 $O(n)$；构造哈希数组另需时间、空间 $O(n)$。}
    {$0\leq n<2^{31}-1$，实际受内存限制；$0,1$ 均不视为素数。}

\subsubsection*{埃氏筛}
按从小到大的顺序枚举未标记的 $p$，它一定为素数，再标记其倍数。可从 $p^2$ 开始，因为小于 $p^2$ 的合数倍数已被更小素因子标记。总复杂度 $O(n\log\log n)$，空间 $O(n)$。若自行实现埃氏筛，应先判断 $p\leq n/p$ 再计算 $p^2$；下方代码采用线性筛。

\subsubsection*{线性筛为何只标记一次}
枚举 $i$，再按升序枚举素数 $p$，标记 $ip$；遇到第一个整除 $i$ 的 $p$ 后停止。这样只使用 $p\leq\operatorname{minp}(i)$ 的组合。

对任意合数 $v$，令 $p=\operatorname{minp}(v)$、$i=v/p$，必有 $p\leq\operatorname{minp}(i)$，所以它会被生成。反过来，每次生成 $ip$ 时 $p$ 恰为其最小质因子，因而同一合数不会由其他组合重复生成。总时间为 $O(n)$。

同时维护欧拉函数：若 $p\mid i$，则 $\varphi(ip)=p\varphi(i)$；否则 $\varphi(ip)=(p-1)\varphi(i)$。两种情况分别对应增加已有质因子的指数、加入一个新的质因子。

莫比乌斯函数同样在线性筛中维护：$p\mid i$ 时 $\mu(ip)=0$，否则 $\mu(ip)=-\mu(i)$，定义与反演见第~\ref{algo:mobius}~节。

\subsubsection*{质因数奇偶哈希}
给每个素数 $p$ 分配一个 64 位随机数 $w_p$，定义
\[
    h(x)=\bigoplus_{p:\,v_p(x)\text{ 为奇数}}w_p.
\]
因为异或两次会抵消，有 $h(ab)=h(a)\mathbin{\oplus}h(b)$。若 $ab$ 为完全平方数，其指纹一定为 $0$；指纹为 $0$ 则仍可能是碰撞，不能作为确定性证明。这里记录的是质因数指数的奇偶性，不是整数的唯一编码。

\subsubsection*{使用教程}
\InlineCode{init(n)} 预处理并覆盖旧表；\InlineCode{primes} 是从 0 开始的升序素数表，\InlineCode{minp[x]}、\InlineCode{phi[x]} 分别为 $x$ 的最小质因子与欧拉函数，\InlineCode{mu[x]}、\InlineCode{pre_mu[x]} 为莫比乌斯函数及其前缀和。当 $2\le x\le n$ 时，\InlineCode{minp[x]==x} 表示素数，反复除以最小质因子可分解 $x$。\InlineCode{prime_xor(seed)} 返回各数的奇偶指纹数组，比较时使用同一随机种子生成的表。
\begin{lstlisting}[style=mtycpp,numbers=none]
Sieve sv;
sv.init(100);
int p=sv.minp[84];
long long val=sv.phi[36];
auto h=sv.prime_xor(20261015);
auto tag=h[12]^h[3];
\end{lstlisting}

\needspace{7\baselineskip}
\subsubsection*{模板代码}
\CodeFile{code/number-theory/sieve.cpp}
